% =====================================================================
%  Monster Tamers — Riepilogo della consegna e dei requisiti funzionali
%  Advanced Software Engineering (Lab) — Università di Pisa
%  Fonti: "Project Kickoff.pdf" (A. Bocci, 07/10/2026)
%         "Monster_Tamers_Backlog.xlsx" (foglio "Project backlog")
%  Compilazione: tectonic -X compile Monster_Tamers_Riepilogo.tex
% =====================================================================
\documentclass[11pt,a4paper]{article}

% ---------------------------------------------------------------------
%  Pacchetti
% ---------------------------------------------------------------------
\usepackage{fontspec}
\usepackage[italian]{babel}
\usepackage[a4paper,margin=2cm,headheight=14pt]{geometry}
\usepackage[table]{xcolor}
\usepackage{tikz}
\usetikzlibrary{arrows.meta,positioning,calc,decorations.pathreplacing,shapes.misc}
\usepackage[most]{tcolorbox}
\usepackage{array,tabularx,booktabs}
\usepackage{enumitem}
\usepackage{titlesec}
\usepackage{fancyhdr}
\usepackage{parskip}
\usepackage{needspace}
\usepackage[hidelinks]{hyperref}

% ---------------------------------------------------------------------
%  Font: tutto il documento in sans (TeX Gyre Heros)
% ---------------------------------------------------------------------
\setmainfont{texgyreheros}[
  Extension      = .otf,
  UprightFont    = *-regular,
  BoldFont       = *-bold,
  ItalicFont     = *-italic,
  BoldItalicFont = *-bolditalic]
\setsansfont{texgyreheros}[
  Extension      = .otf,
  UprightFont    = *-regular,
  BoldFont       = *-bold,
  ItalicFont     = *-italic,
  BoldItalicFont = *-bolditalic]
\setmonofont{lmmono10-regular}[Extension=.otf, Scale=MatchLowercase]

% ---------------------------------------------------------------------
%  Colori
% ---------------------------------------------------------------------
\definecolor{accent}{HTML}{1F4E78}
\definecolor{accentlight}{HTML}{5B9BD5}
\definecolor{prRed}{HTML}{C62828}
\definecolor{prYellow}{HTML}{F9A825}
\definecolor{prGreen}{HTML}{2E7D32}
\definecolor{softgray}{HTML}{F4F6F8}
\definecolor{midgray}{HTML}{6B7280}

% ---------------------------------------------------------------------
%  Titoli, intestazioni, elenchi
% ---------------------------------------------------------------------
\titleformat{\section}{\Large\bfseries\color{accent}}{\thesection}{0.6em}{}[{\color{accentlight}\titlerule[1pt]}]
\titleformat{\subsection}{\large\bfseries\color{accent}}{\thesubsection}{0.6em}{}
\titleformat{\subsubsection}{\normalsize\bfseries\color{accent}}{}{0em}{}
\titlespacing*{\section}{0pt}{1.6em}{0.8em}
\titlespacing*{\subsection}{0pt}{1.2em}{0.4em}

\pagestyle{fancy}
\fancyhf{}
\fancyhead[L]{\small\color{midgray}Monster Tamers --- Riepilogo progetto ASE}
\fancyhead[R]{\small\color{midgray}A.A.\ 2026/27}
\fancyfoot[C]{\small\color{midgray}\thepage}
\renewcommand{\headrulewidth}{0.4pt}
\renewcommand{\headrule}{\hbox to\headwidth{\color{accentlight}\leaders\hrule height \headrulewidth\hfill}}

\setlist{itemsep=2pt,topsep=4pt}
\setlist[itemize]{label={\color{accent}\textbullet}}

% ---------------------------------------------------------------------
%  Macro
% ---------------------------------------------------------------------
\newcommand{\badge}[3]{{\setlength{\fboxsep}{2.2pt}\colorbox{#1}{\textcolor{#2}{\scriptsize\bfseries\strut #3}}}}
\newcommand{\rosso}{\badge{prRed}{white}{PRIMA CONSEGNA}}
\newcommand{\giallo}{\badge{prYellow}{black}{CONSEGNA FINALE}}
\newcommand{\verde}{\badge{prGreen}{white}{OPZIONALE}}
\newcommand{\enc}{\emph{encounter}}
\newcommand{\us}[1]{\textbf{\color{accent}US-#1}}

\newcolumntype{L}[1]{>{\raggedright\arraybackslash}p{#1}}
\newcolumntype{Y}{>{\raggedright\arraybackslash}X}
\newcolumntype{C}[1]{>{\centering\arraybackslash}p{#1}}

% Box generici
\newtcolorbox{infobox}[1][]{enhanced,breakable,colback=accent!5,colframe=accent,
  boxrule=0.8pt,arc=3pt,left=8pt,right=8pt,top=6pt,bottom=6pt,
  fonttitle=\bfseries,coltitle=white,#1}
\newtcolorbox{notebox}[1][]{enhanced,breakable,colback=softgray,colframe=softgray,
  borderline west={3pt}{0pt}{accentlight},arc=0pt,left=10pt,right=8pt,top=5pt,bottom=5pt,#1}
\newtcolorbox{warnbox}[1][]{enhanced,breakable,colback=prRed!5,colframe=prRed!5,
  borderline west={3pt}{0pt}{prRed},arc=0pt,left=10pt,right=8pt,top=5pt,bottom=5pt,#1}
\newtcolorbox{choicebox}[1][]{enhanced,breakable,colback=prYellow!8,colframe=prYellow!80!black,
  boxrule=0.8pt,arc=3pt,left=8pt,right=8pt,top=6pt,bottom=6pt,
  fonttitle=\bfseries,coltitle=black,colbacktitle=prYellow!60,#1}

% Tabella delle user stories
\newenvironment{ustable}{%
  \par\smallskip
  \rowcolors{2}{softgray}{white}%
  \renewcommand{\arraystretch}{1.45}%
  \tabularx{\linewidth}{@{}L{1.45cm}L{3.5cm}Y C{3.0cm}@{}}
  \rowcolor{accent}
  \textcolor{white}{\bfseries\small ID} &
  \textcolor{white}{\bfseries\small Titolo} &
  \textcolor{white}{\bfseries\small Cosa vuole il giocatore e perché} &
  \textcolor{white}{\bfseries\small Quando} \\
}{%
  \endtabularx\par\medskip
}
\newcommand{\usrow}[5]{\us{#1} & \textbf{#2} & \small Voglio #3, \emph{così che} #4. & #5 \\}

% =====================================================================
\begin{document}

% ---------------------------------------------------------------------
%  Frontespizio
% ---------------------------------------------------------------------
\thispagestyle{empty}
\begin{tcolorbox}[enhanced,colback=accent,colframe=accent,arc=4pt,
  left=14pt,right=14pt,top=16pt,bottom=16pt]
  \color{white}
  {\small ADVANCED SOFTWARE ENGINEERING (LAB) --- UNIVERSITÀ DI PISA}\par\medskip
  {\fontsize{30}{34}\selectfont\bfseries Monster Tamers}\par\smallskip
  {\Large Riepilogo della consegna e dei requisiti funzionali}\par\medskip
  {\raggedright\small Docente: Alessandro Bocci\par
  Fonti: \emph{Lab Project Kick-off} (07/10/2026) e \emph{Monster\_Tamers\_Backlog.xlsx}\par}
\end{tcolorbox}

\vspace{4mm}

\begin{infobox}[title={In breve}]
\begin{itemize}[leftmargin=1.2em]
  \item \textbf{Cosa:} un gioco di collezione di mostri (\emph{Monster Tamers}) realizzato come
        \textbf{applicazione distribuita} con backend \textbf{REST}, uno o più \textbf{database}
        e un \textbf{frontend minimale}.
  \item \textbf{Chi:} gruppi da \textbf{3 o 4 persone}, da formare entro il \textbf{1 novembre}.
  \item \textbf{Scadenza finale:} \textbf{15 dicembre}. Discussione orale a \textbf{gennaio/febbraio}.
  \item \textbf{Cosa si consegna:} repository Git (codice + documentazione) e un breve report
        su template fornito dal docente.
  \item \textbf{Requisiti:} 28 user stories --- \textbf{18} per la prima consegna, \textbf{2} in più
        per la consegna finale, \textbf{8} opzionali.
\end{itemize}
\end{infobox}

\vspace{2mm}
\setcounter{tocdepth}{1}
{\small\tableofcontents}

\newpage

% =====================================================================
\section{Organizzazione e scadenze}
% =====================================================================

\subsection{Gruppi}
\begin{itemize}
  \item Gruppi da \textbf{3 o 4 persone}.
  \item Per formarli si usa il \textbf{documento condiviso su Moodle}: seguire le istruzioni del primo
        foglio per cercare compagni e ``iscrivere'' il gruppo.
  \item \textbf{1 novembre}: la modifica del file viene bloccata, quindi il gruppo deve essere iscritto entro quella data.
\end{itemize}

\subsection{Roadmap}

\begin{center}
\begin{tikzpicture}[x=1cm,y=1cm,
  ms/.style={circle,draw=#1,fill=white,line width=1.6pt,minimum size=10pt,inner sep=0pt},
  dt/.style={font=\bfseries\small,text=#1},
  lab/.style={align=center,font=\footnotesize,text width=2.4cm}]
  \draw[-{Stealth[length=3.5mm]},line width=1.2pt,accentlight] (-0.3,0) -- (15.95,0);
  % 1
  \node[ms=accent]   (a) at (0.8,0) {};
  \draw[accent,dashed] (a) -- (0.8,0.9);
  \node[dt=accent,above] at (0.8,0.9) {7 ott};
  \node[lab,above] at (0.8,1.3) {Kickoff del progetto};
  % 2
  \node[ms=midgray]  (b) at (3.15,0) {};
  \draw[midgray,dashed] (b) -- (3.15,-0.9);
  \node[dt=midgray,below] at (3.15,-0.9) {20 ott};
  \node[lab,below] at (3.15,-1.3) {Architettura generale\\\emph{(opzionale)}};
  % 3
  \node[ms=accent]   (c) at (5.5,0) {};
  \draw[accent,dashed] (c) -- (5.5,0.9);
  \node[dt=accent,above] at (5.5,0.9) {21 ott};
  \node[lab,above] at (5.5,1.3) {Pubblicazione documento dei requisiti};
  % 4
  \node[ms=prRed]    (d) at (7.85,0) {};
  \draw[prRed,dashed] (d) -- (7.85,-0.9);
  \node[dt=prRed,below] at (7.85,-0.9) {1 nov};
  \node[lab,below] at (7.85,-1.3) {Scadenza iscrizione gruppi};
  % 5
  \node[ms=midgray]  (e) at (10.2,0) {};
  \draw[midgray,dashed] (e) -- (10.2,0.9);
  \node[dt=midgray,above] at (10.2,0.9) {10 nov};
  \node[lab,above] at (10.2,1.3) {Primo prototipo\\\emph{(opzionale)}};
  % 6
  \node[ms=prRed]    (f) at (12.55,0) {};
  \draw[prRed,dashed] (f) -- (12.55,-0.9);
  \node[dt=prRed,below] at (12.55,-0.9) {15 dic};
  \node[lab,below] at (12.55,-1.3) {\textbf{Consegna finale}};
  % 7
  \node[ms=prRed]    (g) at (14.9,0) {};
  \draw[prRed,dashed] (g) -- (14.9,0.9);
  \node[dt=prRed,above] at (14.9,0.9) {gen/feb};
  \node[lab,above] at (14.9,1.3) {Discussione del progetto};
  % legenda
  \node[font=\scriptsize,anchor=west] at (-0.3,-3.3)
    {\tikz\node[ms=prRed]{};\ obbligatorio \quad
     \tikz\node[ms=midgray]{};\ opzionale (solo feedback) \quad
     \tikz\node[ms=accent]{};\ evento informativo};
\end{tikzpicture}
\end{center}

\begin{center}
\rowcolors{2}{softgray}{white}
\renewcommand{\arraystretch}{1.4}
\begin{tabularx}{\linewidth}{@{}L{2.7cm}L{4.4cm}Y@{}}
  \rowcolor{accent}
  \textcolor{white}{\bfseries Data} & \textcolor{white}{\bfseries Evento} & \textcolor{white}{\bfseries Note} \\
  7 ottobre    & Kickoff                         & Presentazione del progetto. \\
  20 ottobre   & Consegna architettura generale  & \textbf{Opzionale}. Dettagli nella Sezione~\ref{sec:arch}. \\
  21 ottobre   & Pubblicazione requisiti         & Il docente pubblica il documento con i requisiti del progetto. \\
  1 novembre   & Scadenza gruppi                 & Il documento condiviso su Moodle viene bloccato. \\
  10 novembre  & Consegna primo prototipo        & \textbf{Opzionale}. Dettagli nella Sezione~\ref{sec:proto}. \\
  15 dicembre  & \textbf{Consegna finale}        & \textbf{Obbligatoria}. Dettagli nella Sezione~\ref{sec:final}. \\
  gennaio e febbraio & Discussione                 & Orale di gruppo con valutazione individuale (Sezione~\ref{sec:orale}). \\
\end{tabularx}
\end{center}

\subsection{Perché ci sono consegne opzionali}
\begin{notebox}
Le consegne opzionali servono \textbf{solo a ricevere feedback}: \textbf{non hanno voto} e
\textbf{non c'è penalità} se non si consegna. Sono un'occasione per scoprire presto problemi di
progettazione, perché un feedback anticipato costa di solito meno di un refactoring successivo.
Se non si consegna l'architettura non succede nulla, si perde solo il feedback anticipato.
\end{notebox}

In caso di dubbi o blocchi si può scrivere al docente via email e/o andare a ricevimento.

% =====================================================================
\section{Cosa consegnare}
% =====================================================================

\subsection{Consegna 1 --- Architettura generale \badge{midgray}{white}{OPZIONALE · 20 OTT}}\label{sec:arch}

Uno \textbf{schizzo di un'architettura software monolitica} che contenga:
\begin{enumerate}[label=\textbf{\arabic*.}]
  \item i \textbf{componenti principali};
  \item la \textbf{comunicazione con gli utenti} (quali operazioni fa l'utente e verso quale componente);
  \item il \textbf{contenuto del/i database};
  \item le \textbf{user stories complete} (con criteri di accettazione e dipendenze, vedi Sezione~\ref{sec:completare}).
\end{enumerate}

\subsubsection{Formato}
File \textbf{.pdf di slide} (con l'account Unipi sono disponibili gratis Office e Google Docs), possibilmente
\textbf{sotto le 10 slide}, caricato su \textbf{Moodle} (sezione progetto) da \textbf{un solo membro} del gruppo.

\begin{center}
\rowcolors{2}{softgray}{white}
\renewcommand{\arraystretch}{1.4}
\begin{tabularx}{0.92\linewidth}{@{}L{2.6cm}Y@{}}
  \rowcolor{accent}
  \textcolor{white}{\bfseries Slide} & \textcolor{white}{\bfseries Contenuto} \\
  1        & Nome del gruppo e membri. \\
  2        & Rappresentazione delle informazioni del giocatore. \\
  3        & Collezione dei mostri e regole degli \enc. \\
  4 (max 2 slide) & User stories complete. \\
  successive & Disegno dell'architettura, interazioni dell'utente e database. \\
\end{tabularx}
\end{center}

\subsubsection{Linee guida per disegnare l'architettura}
\begin{itemize}
  \item Partire dalle \textbf{user stories} per individuare le operazioni degli utenti.
  \item Definire \textbf{componenti con uno scopo chiaro}.
  \item Ricordare che poi andranno \textbf{implementati i componenti e la API RESTful}.
  \item \textbf{KISS} --- \emph{Keep It Simple and Stupid}.
\end{itemize}

\begin{infobox}[title={Esempio del docente: negozio online (riassunto)}]
L'esempio mostra il livello di dettaglio atteso. Ogni operazione dell'utente è associata a un
componente e a un endpoint REST (descrivibile con \textbf{OpenAPI}); ogni componente che deve
ricordare dati ha il suo database.

\medskip
\rowcolors{2}{white}{accent!4}
\renewcommand{\arraystretch}{1.35}
\begin{tabularx}{\linewidth}{@{}L{2.6cm}L{3.6cm}Y L{1.6cm}@{}}
  \textbf{Componente} & \textbf{Operazioni utente} & \textbf{Endpoint REST} & \textbf{DB} \\ \midrule
  Authentication & login & \texttt{PUT /login?user=\{username\}...} & Users \\
  Catalogue      & consultare gli articoli & \texttt{GET /items}, \texttt{GET /items/\{item\_id\}} & Items \\
  Orders         & creare ordini e pagare, controllare lo stato & \texttt{PUT /orders} + JSON, \texttt{PUT /payment} + JSON, \texttt{GET /orders/\{id\_order\}} & Orders \\
  Payments       & interno (verso un \emph{Bank service} esterno) & --- & --- \\
\end{tabularx}

\medskip
\textbf{Errore tipico evidenziato:} nella prima versione dell'esempio l'operazione ``controllare lo stato
dell'ordine'' non aveva un database degli ordini su cui appoggiarsi. Ogni operazione deve avere i dati che le servono.

\textbf{Database:} i DB disegnati possono essere DBMS diversi, database diversi nello stesso DBMS
oppure tabelle diverse dello stesso DB. La scelta è libera.

\textbf{Record d'esempio:} \emph{Users} (username, password, dati di fatturazione, storico, \dots),
\emph{Items} (nome, prezzo, disponibilità, \dots).
\end{infobox}

\subsection{Consegna 2 --- Primo prototipo \badge{midgray}{white}{OPZIONALE · 10 NOV}}\label{sec:proto}

Un \textbf{repository GitHub} contenente:
\begin{itemize}
  \item \textbf{descrizione generale} del progetto;
  \item \textbf{codice};
  \item \textbf{quick start}: istruzioni di installazione ed esecuzione;
  \item \textbf{API completa e relativi test}.
\end{itemize}
Come preparare quick start, API e test verrà spiegato nei prossimi laboratori.
Il prototipo copre le user stories \rosso.

\subsection{Consegna finale \badge{prRed}{white}{OBBLIGATORIA · 15 DIC}}\label{sec:final}
\begin{itemize}
  \item \textbf{Repository GitHub} strutturato come per il prototipo (codice e documentazione).
  \item \textbf{Report} che descrive il progetto, scritto sul \textbf{template fornito dal docente}.
\end{itemize}
La differenza tra prototipo e consegna finale è l'aggiunta delle funzionalità (requisiti) che verranno
trattate più avanti nel corso: in particolare le user stories \giallo.

\subsection{Discussione del progetto (gennaio/febbraio)}\label{sec:orale}
\begin{itemize}
  \item La consegna finale riceve \textbf{lo stesso voto per tutti i membri} del gruppo.
  \item Nel periodo d'esame (gennaio e febbraio) il docente proporrà alcune date per incontrare il gruppo.
  \item Durante la discussione il docente fa domande sulla realizzazione del progetto per valutare il
        \textbf{coinvolgimento individuale}.
  \item La discussione \textbf{aggiunge o toglie punti} al voto della consegna, \textbf{individualmente}.
\end{itemize}
\begin{warnbox}
Se uno studente non sa rispondere e non dimostra un coinvolgimento sufficiente nel progetto, il suo
voto finale è \textbf{``Failed''} e dovrà sostenere l'esame \textbf{senza il progetto}.
\end{warnbox}

% =====================================================================
\needspace{9cm}
\section{Requisiti tecnici del software}
% =====================================================================

\begin{tcbraster}[raster columns=3,raster equal height,raster column skip=3mm,raster row skip=3mm]
  \begin{tcolorbox}[colback=accent!5,colframe=accent!70,title=Applicazione distribuita,fonttitle=\bfseries\small,arc=3pt]
    \small Il sistema deve essere un'applicazione distribuita.
  \end{tcolorbox}
  \begin{tcolorbox}[colback=accent!5,colframe=accent!70,title=Interfaccia REST,fonttitle=\bfseries\small,arc=3pt]
    \small Tutte le funzionalità devono essere utilizzabili tramite un'interfaccia RESTful.
  \end{tcolorbox}
  \begin{tcolorbox}[colback=accent!5,colframe=accent!70,title=Uno o più database,fonttitle=\bfseries\small,arc=3pt]
    \small Persistenza dei dati, incluso lo storico degli \enc.
  \end{tcolorbox}
  \begin{tcolorbox}[colback=accent!5,colframe=accent!70,title=Meccanismi di sicurezza,fonttitle=\bfseries\small,arc=3pt]
    \small Protezione di account e dati di gioco.
  \end{tcolorbox}
  \begin{tcolorbox}[colback=accent!5,colframe=accent!70,title=Testing,fonttitle=\bfseries\small,arc=3pt]
    \small Test unitari e test di performance.
  \end{tcolorbox}
  \begin{tcolorbox}[colback=accent!5,colframe=accent!70,title=Analisi di sicurezza,fonttitle=\bfseries\small,arc=3pt]
    \small Analisi della sicurezza del sistema.
  \end{tcolorbox}
\end{tcbraster}

\subsection{Backend}
Il focus principale del progetto è il \textbf{backend}. Il servizio deve gestire:
\begin{itemize}
  \item la \textbf{collezione globale} dei mostri e le \textbf{collezioni personali} dei giocatori;
  \item la \textbf{selezione del team};
  \item gli \textbf{\enc};
  \item lo \textbf{storico degli \enc}, con \textbf{memorizzazione permanente}.
\end{itemize}
Tutto deve essere eseguibile tramite l'\textbf{interfaccia RESTful}.

\subsection{Frontend}
\begin{itemize}
  \item Richiesto un \textbf{frontend minimale} con una UI di base, usato come sistema di debug,
        che chiama la \textbf{REST API del backend} per le operazioni principali.
  \item L'aspetto grafico è libero: conta come viene usato e come interagisce con il backend.
\end{itemize}
\begin{warnbox}
\textbf{Nessuna logica di business rilevante nel frontend.} Le regole del gioco stanno nel backend.
\end{warnbox}

\subsection{Scelte libere}
Linguaggi di programmazione, database e altri aspetti non specificati sono a scelta del gruppo.
Le funzionalità da implementare sono quelle descritte nelle user stories (Sezione~\ref{sec:us}).

% =====================================================================
\needspace{10cm}
\section{Il gioco: Monster Tamers}
% =====================================================================

\subsection{Il ciclo di gioco (game loop)}

Il ciclo minimo richiesto è \textbf{Prepara $\rightarrow$ Affronta $\rightarrow$ Esito $\rightarrow$ Progresso}.
Si può estendere a piacere.

\begin{center}
\begin{tikzpicture}[
  step/.style={rounded corners=4pt,draw=#1!80!black,fill=#1!12,line width=0.8pt,
               text width=2.1cm,align=center,minimum height=1.4cm,font=\scriptsize,inner sep=4pt},
  arr/.style={-{Stealth[length=2.6mm]},line width=0.9pt,color=black!55},
  brace/.style={decorate,decoration={brace,amplitude=5pt},line width=0.8pt,color=accent},
  phase/.style={font=\small\bfseries,text=accent,above=6pt}]
  \node[step=accent]   (col)  at (0,0)       {\textbf{1. Collezione}\\raccogli e possiedi mostri};
  \node[step=accent]   (team) at (2.7,0)     {\textbf{2. Team}\\forma un team con i tuoi mostri};
  \node[step=prRed]    (enc)  at (5.4,0)     {\textbf{3. Encounter}\\affronta una sfida con il team};
  \node[step=prGreen]  (win)  at (8.4,1.0)   {\textbf{4a. Successo}\\encounter vinto};
  \node[step=prRed]    (lose) at (8.4,-1.0)  {\textbf{4b. Fallimento}\\encounter perso};
  \node[step=violet]   (hist) at (11.4,0)    {\textbf{5. Storico}\\rivedi encounter ed esiti};
  \node[step=prYellow,dashed] (rew) at (14.1,0) {\textbf{Ricompensa}\\\emph{(opzionale)}\\nuovi mostri};

  \draw[arr] (col) -- (team);
  \draw[arr] (team) -- (enc);
  \draw[arr] (enc.east) -- ++(0.3,0) |- (win.west);
  \draw[arr] (enc.east) -- ++(0.3,0) |- (lose.west);
  \draw[arr] (win.east)  -- ++(0.3,0) |- (hist.west);
  \draw[arr] (lose.east) -- ++(0.3,0) |- (hist.west);
  \draw[arr,dashed] (hist) -- (rew);
  \draw[arr,dashed] (rew.south) -- ++(0,-1.75) -| node[pos=0.25,below,font=\scriptsize,text=black!60]
       {nuovi mostri entrano nella collezione: il ciclo ricomincia} (col.south);

  \draw[brace] ($(col.west)+(0,2.0)$)  -- node[phase]{Prepara}   ($(team.east)+(0,2.0)$);
  \draw[brace] ($(enc.west)+(0,2.0)$)  -- node[phase]{Affronta}  ($(enc.east)+(0,2.0)$);
  \draw[brace] ($(win.west)+(0,1.0)$)  -- node[phase]{Esito}     ($(win.east)+(0,1.0)$);
  \draw[brace] ($(hist.west)+(0,2.0)$) -- node[phase]{Progresso} ($(rew.east)+(0,2.0)$);
\end{tikzpicture}
\end{center}

\subsection{Giocatori}
Come rappresentare un giocatore è a scelta del gruppo. \textbf{Minimo obbligatorio:}
\begin{itemize}
  \item delle \textbf{credenziali};
  \item lo \textbf{storico degli \enc}.
\end{itemize}
Si possono aggiungere altri dati (immagine, amici, \dots) e si possono tenere separati
\textbf{profilo} e \textbf{account} del giocatore.

\subsection{Mostri}
\begin{itemize}
  \item Ogni gruppo sceglie la \textbf{collezione di mostri} del proprio progetto.
  \item Ogni mostro deve avere \textbf{almeno}:
    \begin{itemize}
      \item un'\textbf{immagine} che lo rappresenta;
      \item \textbf{due caratteristiche} usate nel gameplay degli \enc{} (se ne possono aggiungere altre).
    \end{itemize}
  \item Le caratteristiche \textbf{influenzano gli \enc}.
  \item Ogni giocatore possiede una \textbf{collezione personale}.
  \item Durante il gioco si possono \textbf{ottenere nuovi mostri}.
\end{itemize}
\begin{notebox}
\textbf{Solo come esempio}, nelle slide i mostri hanno: un elemento (fuoco, acqua, terra, oscurità),
le statistiche \emph{Power}, \emph{Speed}, \emph{Defense}, \emph{Magic} (da 1 a 5 stelle) e un'abilità speciale
(es.\ ``riduce il danno subito del 25\%''). Non è un vincolo.
\end{notebox}

\subsection{Team}
Prima di un \enc{} il giocatore prepara un team con i mostri della propria collezione. Il team deve:
\begin{itemize}
  \item avere una \textbf{dimensione limitata};
  \item \textbf{rispettare le regole del gioco};
  \item essere composto \textbf{solo da mostri posseduti}.
\end{itemize}
Le regole precise dei team sono a scelta del gruppo.

\subsection{Encounter}
Un \enc{} è una sfida tra un team e un \textbf{``mostro selvatico''}.
\begin{itemize}
  \item \enc{} diversi possono avere caratteristiche diverse;
  \item il risultato \textbf{dipende dal team scelto};
  \item un \enc{} può \textbf{riuscire o fallire};
  \item alcuni esiti possono \textbf{dare un nuovo mostro}.
\end{itemize}
Un \enc{} \textbf{non deve per forza} essere una battaglia a turni.

\subsection{Regole del gioco}
Le regole sono in gran parte libere. Ad esempio vanno definiti:
l'obiettivo principale del gioco, come si amplia la collezione, come funzionano gli \enc,
se sono ammessi duplicati nelle collezioni personali, ecc.

\emph{Esempio dalle slide:} \enc{} basati su combattimenti, con l'obiettivo di battere un boss finale
dopo diverse battaglie di difficoltà crescente.

\begin{choicebox}[title={Decisioni lasciate al gruppo}]
\begin{itemize}[leftmargin=1.2em]
  \item Dati del giocatore oltre al minimo (credenziali + storico); profilo e account separati o no.
  \item Collezione dei mostri: quali mostri, quali caratteristiche (almeno immagine + 2 caratteristiche).
  \item Regole dei team: dimensione massima e vincoli.
  \item Meccanica degli \enc: come si calcola l'esito (anche non a turni), quando si ottiene un mostro.
  \item Obiettivo del gioco, crescita della collezione, gestione dei duplicati.
  \item Funzionalità extra (user stories verdi o idee nuove).
  \item Linguaggi, database, aspetto grafico del frontend.
\end{itemize}
\end{choicebox}

% =====================================================================
\section{Requisiti funzionali (user stories)}\label{sec:us}
% =====================================================================

\subsection{Come leggere le priorità}

\begin{center}
\renewcommand{\arraystretch}{1.5}
\begin{tabularx}{0.95\linewidth}{@{}C{3.2cm}Y C{1.6cm}@{}}
  \toprule
  \textbf{Colore} & \textbf{Significato} & \textbf{N.\ US} \\ \midrule
  \rosso   & Funzionalità core, da implementare già nel \textbf{primo prototipo} (e quindi presenti anche nella consegna finale). & \textbf{18} \\
  \giallo  & Funzionalità da aggiungere per la \textbf{consegna finale} (legate ad argomenti trattati più avanti nel corso). & \textbf{2} \\
  \verde   & \textbf{Facoltative}: proposte dagli studenti, si possono implementare oppure no. & \textbf{8} \\ \bottomrule
\end{tabularx}
\end{center}

\begin{notebox}
\begin{itemize}[leftmargin=1.2em]
  \item La priorità indica il \textbf{prodotto minimo funzionante per il progetto}, non per un prodotto reale.
  \item Le user stories sono volutamente \textbf{generiche} per lasciare libertà ai gruppi (ad esempio, non
        esiste una user story su come implementare gli \enc). Quelle specifiche sono tutte verdi.
  \item Si possono aggiungere \textbf{funzionalità non presenti} nelle user stories.
  \item In tutte le user stories l'attore è \textbf{il giocatore} (``Come giocatore, voglio\dots'').
\end{itemize}
\end{notebox}

{\footnotesize\color{midgray} Nota: il file Excel contiene anche un foglio ``Realistic backlog'' con priorità P1--P3
pensate per un gioco reale. Non riguardano il progetto e sono state omesse; da quel foglio sono ripresi solo
i titoli brevi delle user stories.\par}

\subsection{Da completare a cura del gruppo}\label{sec:completare}
Il file del docente riporta solo le user stories con la loro priorità (criteri di accettazione e dipendenze
sono stati omessi di proposito). Il docente chiede di:
\begin{itemize}
  \item definire i \textbf{criteri di accettazione} almeno per le user stories ad \textbf{alta e media priorità};
  \item definire le \textbf{dipendenze} tra le user stories.
\end{itemize}

\subsection{Account e progressi}
\begin{ustable}
  \usrow{01}{Creare un account}{creare un account}{i miei progressi di gioco siano associati a me}{\rosso}
  \usrow{02}{Accedere (login)}{effettuare il login}{io possa accedere al mio gioco e continuare a giocare}{\rosso}
  \usrow{03}{Ritrovare i progressi salvati}{ritrovare i miei progressi quando torno}{io possa riprendere da dove avevo lasciato}{\rosso}
  \usrow{04}{Personalizzare il profilo}{scegliere o modificare il mio nome da giocatore e il mio avatar}{io possa personalizzare la mia identità}{\verde}
  \usrow{05}{Proteggere i dati dell'account}{che i dati del mio account siano protetti}{nessun altro possa accedervi o modificarli senza permesso}{\giallo}
\end{ustable}

\subsection{Collezione di mostri}
\begin{ustable}
  \usrow{06}{Ricevere una collezione iniziale}{ricevere un set iniziale di mostri quando inizio a giocare}{io possa formare il mio primo team}{\rosso}
  \usrow{07}{Scegliere il primo mostro}{scegliere il mio primo mostro}{io possa personalizzare la mia strategia iniziale}{\verde}
  \usrow{08}{Scegliere una fazione}{scegliere una fazione (o un'affiliazione iniziale equivalente)}{la mia scelta iniziale possa influenzare l'esperienza di gioco}{\verde}
  \usrow{09}{Vedere la collezione personale}{vedere i mostri che possiedo}{io sappia quali mostri posso usare}{\rosso}
  \usrow{10}{Sfogliare la collezione globale}{sfogliare tutti i mostri disponibili nel gioco}{io capisca cosa posso incontrare od ottenere}{\rosso}
  \usrow{11}{Vedere i dettagli di un mostro}{vedere l'immagine e le caratteristiche di un mostro}{io possa valutare quanto può essermi utile negli \enc}{\rosso}
  \usrow{12}{Ottenere un mostro guadagnato}{che un mostro guadagnato o catturato in un \enc{} venga aggiunto alla mia collezione}{le mie scelte future si amplino}{\rosso}
  \usrow{13}{Sviluppare i mostri}{migliorare o personalizzare i miei mostri}{io possa sviluppare strategie diverse nel tempo}{\verde}
\end{ustable}

\subsection{Team}
\begin{ustable}
  \usrow{14}{Creare un team}{creare un team usando i mostri della mia collezione}{io possa prepararmi a un \enc}{\rosso}
  \usrow{15}{Modificare un team}{sostituire o rimuovere mostri dal mio team}{io possa adattare la strategia prima di un \enc}{\rosso}
  \usrow{16}{Validare un team}{sapere se il team selezionato rispetta le regole del gioco}{io non inizi con un team non valido}{\rosso}
\end{ustable}

\subsection{Encounter}
\begin{ustable}
  \usrow{17}{Sfogliare gli encounter}{vedere gli \enc{} disponibili e le loro caratteristiche}{io possa scegliere una sfida adatta}{\rosso}
  \usrow{18}{Avviare un encounter}{avviare un \enc{} con il team selezionato}{io possa affrontare una sfida}{\rosso}
  \usrow{19}{Risolvere un encounter}{che l'\enc{} venga risolto secondo le regole del gioco}{io sappia se ho avuto successo o ho fallito}{\rosso}
  \usrow{20}{Dettagli del risultato}{vedere i dettagli del risultato di un \enc}{io capisca l'esito e il team coinvolto}{\rosso}
  \usrow{21}{Abbandonare un encounter}{abbandonare o fuggire da un \enc{} quando le regole lo consentono}{io possa interrompere una sfida che non voglio più continuare}{\verde}
  \usrow{22}{Scegliere le azioni}{scegliere azioni come attaccare o tentare di catturare un mostro durante un \enc}{le mie decisioni possano influenzarne l'esito}{\verde}
\end{ustable}

\subsection{Storico e integrità}
\begin{ustable}
  \usrow{23}{Vedere lo storico}{vedere i miei \enc{} precedenti}{io possa rivedere i miei progressi e i risultati}{\rosso}
  \usrow{24}{Dettagli di un encounter passato}{vedere i dettagli di un \enc{} passato}{io possa imparare dalla scelta del team e dal risultato}{\rosso}
  \usrow{25}{Proteggere i salvataggi}{che collezione, storico e progressi siano protetti da manomissioni o perdite}{io possa fidarmi del mio salvataggio}{\giallo}
\end{ustable}

\subsection{Social}
\begin{ustable}
  \usrow{26}{Gestire gli amici}{aggiungere e gestire amici}{io possa tenere traccia delle persone che conosco nel gioco}{\verde}
  \usrow{27}{Sfidare un amico}{sfidare un amico}{io possa giocare con altri in una modalità aggiuntiva}{\verde}
\end{ustable}

\subsection{Interfaccia del giocatore}
\begin{ustable}
  \usrow{28}{Usare un'interfaccia semplice}{eseguire le azioni principali del gioco tramite un'interfaccia semplice}{io possa giocare senza usare strumenti tecnici}{\rosso}
\end{ustable}

\subsection{Vista per consegna}

\begin{tcbraster}[raster columns=3,raster equal height,raster column skip=3mm]
  \begin{tcolorbox}[colback=prRed!5,colframe=prRed,title={Primo prototipo (18)},fonttitle=\bfseries\small,arc=3pt,left=5pt,right=5pt]
    \footnotesize
    \us{01} Creare account\\
    \us{02} Login\\
    \us{03} Ritrovare progressi\\
    \us{06} Collezione iniziale\\
    \us{09} Collezione personale\\
    \us{10} Collezione globale\\
    \us{11} Dettagli mostro\\
    \us{12} Ottenere mostro\\
    \us{14} Creare team\\
    \us{15} Modificare team\\
    \us{16} Validare team\\
    \us{17} Sfogliare encounter\\
    \us{18} Avviare encounter\\
    \us{19} Risolvere encounter\\
    \us{20} Dettagli risultato\\
    \us{23} Storico\\
    \us{24} Dettagli encounter passato\\
    \us{28} Interfaccia semplice
  \end{tcolorbox}
  \begin{tcolorbox}[colback=prYellow!10,colframe=prYellow!90!black,coltitle=black,title={+ Consegna finale (2)},fonttitle=\bfseries\small,arc=3pt,left=5pt,right=5pt]
    \footnotesize
    \us{05} Protezione dati account\\
    \us{25} Protezione salvataggi da manomissioni o perdite

    \medskip
    \textit{Tutte e due riguardano la sicurezza e l'integrità dei dati, in linea con i requisiti tecnici
    ``meccanismi di sicurezza'' e ``analisi di sicurezza''.}

    \medskip
    \textit{La consegna finale comprende quindi tutte le 18 US rosse più queste 2.}
  \end{tcolorbox}
  \begin{tcolorbox}[colback=prGreen!6,colframe=prGreen,title={Opzionali (8)},fonttitle=\bfseries\small,arc=3pt,left=5pt,right=5pt]
    \footnotesize
    \us{04} Nome e avatar\\
    \us{07} Scegliere primo mostro\\
    \us{08} Scegliere fazione\\
    \us{13} Sviluppare mostri\\
    \us{21} Abbandonare encounter\\
    \us{22} Scegliere azioni\\
    \us{26} Gestire amici\\
    \us{27} Sfidare amico
  \end{tcolorbox}
\end{tcbraster}

\end{document}
